%=============================================================================
% MÓDULO 05: METODOLOGIA (Expansão Final)
%=============================================================================
% Frontiers in Public Health — Alinhada ao PROBAST, reproduzível
%   5.1 Desenho do Estudo
%   5.2 Fonte de Dados
%   5.3 Engenharia de Características
%   5.4 Modelos
%   5.5 Métricas de Justiça
%   5.6 Análise Interseccional
%   5.7 Implementação
%=============================================================================

\section{Metodologia}

\subsection{Desenho do Estudo}

Este estudo segue as diretrizes do Prediction model Risk Of Bias ASsessment 
Tool (PROBAST) \cite{wolff2019} para avaliar viés em modelos de predição. 
Conduzimos uma análise retrospectiva de um conjunto de dados de diabetes 
publicamente disponível, aplicando quatro algoritmos de aprendizado de 
máquina e avaliando o desempenho em subgrupos demográficos definidos por 
gênero e etnia.

O arcabouço PROBAST fornece uma abordagem estruturada para avaliar o risco 
de viés em estudos de modelos de predição. Ele avalia o viés em quatro 
domínios: participantes, preditores, desfecho e análise. Nosso estudo 
aborda cada domínio sistematicamente, conforme detalhado abaixo.

\subsubsection{Domínio 1: Participantes}

O Banco de Dados de Diabetes dos Índios Pima inclui 768 mulheres adultas 
(idade ≥ 21) de herança indígena Pima, coletado pelo Instituto Nacional 
de Diabetes e Doenças Digestivas e Renais (NIDDK). O conjunto de dados 
representa uma população de alto risco com prevalência de 34.9\%, tornando-o 
adequado para avaliar desempenho e justiça de modelos de predição.

O conjunto de dados original foi coletado entre 1986 e 1991 como parte de 
um estudo epidemiológico maior sobre diabetes entre os Índios Pima. Embora 
a população específica do conjunto de dados limite a generalizabilidade, 
ele fornece uma coorte bem caracterizada com medições clínicas abrangentes 
que permitem uma avaliação completa do desempenho e justiça do modelo.

\subsubsection{Domínio 2: Preditores}

Oito preditores clínicos foram selecionados com base em associações 
estabelecidas com risco de diabetes e disponibilidade no conjunto de dados:

\begin{itemize}[leftmargin=*]
    \item \textbf{Glicose}: Concentração de glicose plasmática em 2 horas 
          em teste oral de tolerância à glicose (mg/dL). Glicose em jejum 
          ≥ 126 mg/dL indica diabetes; glicose em 2 horas ≥ 200 mg/dL 
          indica diabetes.
    \item \textbf{IMC}: Índice de massa corporal calculado como peso (kg) 
          dividido por altura (m)$^2$. IMC ≥ 30 indica obesidade, um 
          fator de risco importante para diabetes tipo 2.
    \item \textbf{Idade}: Idade em anos no momento do exame. O risco de 
          diabetes aumenta com a idade, particularmente após os 45 anos.
    \item \textbf{Insulina}: Nível sérico de insulina em 2 horas (μU/mL). 
          Níveis elevados de insulina podem indicar resistência à insulina, 
          um precursor do diabetes.
    \item \textbf{Espessura da Pele}: Espessura da dobra cutânea do tríceps 
          (mm), uma medida de gordura subcutânea que se correlaciona com a 
          composição corporal e risco metabólico.
    \item \textbf{Gestações}: Número de gestações, um fator de risco para 
          diabetes gestacional e subsequente diabetes tipo 2.
    \item \textbf{Pressão Arterial}: Pressão arterial diastólica (mm Hg). 
          A hipertensão é frequentemente associada a diabetes e síndrome 
          metabólica.
    \item \textbf{Glicose\_IMC\_Idade}: Termo de interação de três vias 
          capturando efeitos sinérgicos de fatores metabólicos (glicose × 
          IMC × idade / 1000), projetado para capturar a interação complexa 
          entre esses fatores de risco chave.
\end{itemize}

A seleção desses preditores segue o conhecimento clínico estabelecido sobre 
fatores de risco para diabetes e garante que o modelo incorpore tanto 
características tradicionais quanto baseadas em interação. O termo de 
interação foi incluído para capturar relações não lineares entre preditores 
chave que podem ser clinicamente relevantes.

\subsubsection{Domínio 3: Desfecho}

O desfecho binário foi o status de diabetes diagnosticado de acordo com 
os critérios da OMS: glicose plasmática em 2 horas ≥ 200 mg/dL durante 
teste oral de tolerância à glicose. Esse critério diagnóstico é amplamente 
utilizado na prática clínica e fornece uma definição padronizada para 
classificação de casos.

\subsubsection{Domínio 4: Análise}

O desenvolvimento do modelo seguiu um protocolo analítico rigoroso:

\begin{enumerate}[leftmargin=*]
    \item \textbf{Divisão de dados}: Divisão treino-teste 80/20 estratificada 
          por desfecho para manter a distribuição de classes entre as 
          partições. Isso garante que os conjuntos de treinamento e teste 
          tenham prevalência similar de diabetes, permitindo avaliação justa.
    \item \textbf{Padronização de características}: StandardScaler aplicado 
          para normalizar características com média zero e variância unitária, 
          essencial para algoritmos baseados em distância como SVM.
    \item \textbf{Validação cruzada}: Validação cruzada estratificada em 5 
          dobras no conjunto de treinamento para avaliar estabilidade do 
          modelo e reduzir risco de sobreajuste. A validação cruzada fornece 
          uma estimativa mais robusta de desempenho do que divisões únicas 
          treino-teste.
    \item \textbf{Desequilíbrio de classes}: Pesos de classe balanceados 
          aplicados para abordar a prevalência de 34.9\%, impedindo que os 
          modelos sejam enviesados em direção à classe majoritária.
    \item \textbf{Hiperparâmetros}: Parâmetros padrão do scikit-learn 
          utilizados para reprodutibilidade e evitar sobreajuste por meio 
          de ajuste excessivo. Embora a otimização de hiperparâmetros possa 
          melhorar o desempenho, também pode introduzir viés e reduzir a 
          reprodutibilidade.
\end{enumerate}

\subsection{Pré-processamento de Dados}

O pré-processamento de dados abordou problemas comuns de qualidade em 
conjuntos de dados clínicos:

\begin{enumerate}[leftmargin=*]
    \item \textbf{Valores ausentes}: Imputados usando valores medianos para 
          características com valores fisiológicos zero (glicose, IMC, pressão 
          arterial, insulina, espessura da pele). O método de imputação 
          mediana foi escolhido por sua robustez a outliers e preservação 
          da distribuição de dados.
    \item \textbf{Outliers}: Mantidos como anomalias clínicas, sem aplicação 
          de truncamento. Valores extremos em dados clínicos podem representar 
          casos genuinamente raros mas clinicamente importantes, e removê-los 
          pode enviesar o modelo.
    \item \textbf{Características de interação}: Criada glicose\_IMC\_idade 
          para capturar efeitos sinérgicos entre preditores metabólicos chave. 
          Esse termo de interação foi projetado com base em conhecimento 
          clínico sobre como esses fatores interagem para influenciar o 
          risco de diabetes.
\end{enumerate}

\subsection{Modelos de Aprendizado de Máquina}

Quatro algoritmos com paradigmas de aprendizado distintos foram selecionados 
para avaliar se o viés é específico do algoritmo ou sistêmico:

\begin{enumerate}[leftmargin=*]
    \item \textbf{Regressão Logística} \cite{cox1958}: Modelo linear para 
          classificação binária, interpretável e amplamente utilizado em 
          contextos clínicos. A regressão logística fornece uma linha de 
          base para separabilidade linear e oferece interpretação transparente 
          de coeficientes.
    \item \textbf{Random Forest} \cite{breiman2001}: Ensemble de árvores 
          de decisão com bagging, robusto a sobreajuste. Random forests 
          capturam relações não lineares e interações sem engenharia de 
          características explícita.
    \item \textbf{Gradient Boosting} \cite{friedman2001}: Ensemble sequencial 
          com otimização por gradiente descendente, tipicamente de alto 
          desempenho. Gradient boosting constrói árvores sequencialmente, 
          com cada árvore corrigindo erros de iterações anteriores.
    \item \textbf{Support Vector Machine} (SVM) \cite{vapnik1995}: 
          Classificador baseado em kernel com maximização de margem, eficaz 
          em espaços de alta dimensionalidade. SVMs são particularmente 
          eficazes quando o número de características é grande em relação 
          ao número de amostras.
\end{enumerate}

Esses algoritmos representam diferentes famílias de métodos de aprendizado 
de máquina: modelos lineares (regressão logística), ensembles baseados em 
árvores (random forest, gradient boosting) e métodos de kernel (SVM). Ao 
avaliar múltiplos paradigmas algorítmicos, podemos determinar se os vieses 
observados são específicos de certos tipos de modelos ou representam 
questões sistêmicas mais amplas.

\subsection{Métricas de Justiça}

Empregamos quatro métricas complementares de justiça para fornecer uma 
avaliação abrangente do viés algorítmico:

\subsubsection{Odds Equalizados}

A diferença de Odds Equalizados (EOD) mede a discrepância nas taxas de 
erro entre os grupos \cite{hardt2016}:

\begin{equation}
\text{EOD} = |(\text{TPR}_a - \text{TPR}_b)| + |(\text{FPR}_a - \text{FPR}_b)|
\label{eq:eod}
\end{equation}

\noindent onde TPR é a taxa de verdadeiros positivos e FPR é a taxa de 
falsos positivos para os grupos $a$ e $b$. O EOD captura disparidades 
tanto em diagnósticos perdidos (falsos negativos) quanto em intervenções 
desnecessárias (falsos positivos), fornecendo uma medida holística de 
justiça.

\subsubsection{Diferença de Paridade Demográfica}

A Diferença de Paridade Demográfica (DPD) quantifica a diferença absoluta 
nas taxas de predição positiva entre os grupos \cite{dwork2012}:

\begin{equation}
\text{DPD} = |P(\hat{Y}=1 \mid A=a) - P(\hat{Y}=1 \mid A=b)|
\label{eq:dpd}
\end{equation}

O DPD mede se o modelo recomenda intervenções igualmente entre os grupos, 
independente da prevalência real da doença. Embora o DPD não leve em conta 
diferenças legítimas nas taxas básicas, ele fornece uma medida importante 
de acesso equitativo à saúde assistida por IA.

\subsubsection{Disparidade de Taxa de Verdadeiros Positivos}

A disparidade de TPR mede a diferença na sensibilidade entre os grupos, 
capturando se o modelo é igualmente eficaz na identificação de casos 
positivos em todas as populações:

\begin{equation}
\text{Disparidade de TPR} = \text{TPR}_a - \text{TPR}_b
\label{eq:tpr_disparity}
\end{equation}

Essa métrica é particularmente importante em contextos de triagem clínica, 
onde diagnósticos perdidos (falsos negativos) podem ter consequências 
sérias de saúde. Uma disparidade positiva de TPR indica que o grupo $a$ 
tem sensibilidade mais alta que o grupo $b$, enquanto um valor negativo 
indica o contrário.

\subsubsection{Disparidade de Taxa de Falsos Positivos}

A disparidade de FPR mede a diferença nas taxas de alarme falso entre 
os grupos, capturando se o modelo produz mais intervenções desnecessárias 
para certas populações:

\begin{equation}
\text{Disparidade de FPR} = \text{FPR}_a - \text{FPR}_b
\label{eq:fpr_disparity}
\end{equation}

Alta disparidade de FPR indica que um grupo experimenta mais alarmes 
falsos, levando a ansiedade desnecessária, testes adicionais e potencial 
sobretratamento. Na triagem de diabetes, falsos positivos podem resultar 
em modificações de estilo de vida desnecessárias, medicação e carga 
psicológica.

\subsection{Análise Interseccional}

Estendemos a análise de justiça de atributo único para examinar viés na 
intersecção de gênero e etnia. Essa abordagem, informada pela teoria da 
interseccionalidade \cite{crenshaw1989}, examina se indivíduos na 
intersecção de múltiplas identidades marginalizadas experimentam 
desvantagem amplificada que não pode ser capturada analisando cada atributo 
separadamente.

A análise interseccional envolve os seguintes passos:

\begin{enumerate}[leftmargin=*]
    \item \textbf{Definição do grupo}: Criar subgrupos definidos por todas 
          as combinações de gênero (masculino/feminino) e etnia 
          (majoritário/minoritário), produzindo quatro grupos 
          interseccionais: majoritário--feminino, majoritário--masculino, 
          minoritário--feminino e minoritário--masculino.
    \item \textbf{Avaliação de desempenho}: Calcular precisão, F1-score e 
          TPR para cada grupo interseccional, fornecendo uma avaliação 
          abrangente do desempenho do modelo em todo o espaço demográfico.
    \item \textbf{Quantificação de disparidades}: Comparar métricas de 
          desempenho em todos os grupos interseccionais para identificar 
          os subgrupos mais e menos favorecidos. Essa comparação revela 
          se combinações específicas de atributos demográficos levam a 
          desvantagem amplificada.
    \item \textbf{Avaliação do efeito composto}: Determinar se as 
          disparidades na intersecção excedem a soma das disparidades de 
          atributo único, indicando viés composto. Essa avaliação é crítica 
          para entender se os efeitos interseccionais criam modos únicos 
          de discriminação.
\end{enumerate}

\subsection{Implementação}

Todos os experimentos foram implementados em Python usando scikit-learn 
1.9.0, com os seguintes pacotes principais: numpy 2.5.1, pandas 3.0.3, 
scipy 1.18.0 e matplotlib para visualização. O pipeline experimental foi 
projetado para reprodutibilidade, com sementes aleatórias fixas e operações 
determinísticas.

A análise foi conduzida em uma estação de trabalho com 8GB de RAM e 
processador multi-core, com tempo de execução total de aproximadamente 15 
minutos para todos os experimentos. O código está publicamente disponível 
em \url{https://github.com/MarceloClaro/opencode-ecosystem-core}, permitindo 
verificação independente e extensão de nossas descobertas.

Detalhes-chave de implementação incluem:

\begin{itemize}[leftmargin=*]
    \item Semente aleatória definida como 42 para reprodutibilidade em 
          todos os experimentos
    \item Amostragem estratificada utilizada para divisão treino-teste e 
          validação cruzada para manter a distribuição de classes
    \item StandardScaler ajustado apenas nos dados de treinamento e aplicado 
          tanto nos dados de treinamento quanto de teste para evitar vazamento 
          de dados
    \item Pesos de classe definidos como ``balanceados'' para abordar 
          desequilíbrio de classes sem sobreamostragem sintética
    \item Todos os hiperparâmetros definidos como padrão do scikit-learn 
          para garantir reprodutibilidade e evitar sobreajuste por meio 
          de ajuste
\end{itemize}
